%% rebuttal_preamble.tex  版式与命令定义，一般不需要改。
%% 由主文件读入。编译 xelatex（文档类 ctexart，fontset=none）。

%% ---- 页面 ------------------------------------------------------------------
\usepackage[a4paper,margin=14mm,top=13mm,bottom=15mm]{geometry}
\usepackage{xcolor}
\usepackage{amsmath,amssymb}
\usepackage[most]{tcolorbox}
\usepackage{booktabs,array,tabularx,enumitem,newunicodechar,adjustbox}
\usepackage[hidelinks,bookmarksnumbered]{hyperref}
\setlist{nosep,leftmargin=1.6em}
\setlength{\parindent}{0pt}
\setlength{\parskip}{2pt}
\linespread{1.08}
\ctexset{
  section={format=\large\bfseries\raggedright,beforeskip=8pt,afterskip=4pt},
  subsection={format=\normalsize\bfseries\raggedright,beforeskip=6pt,afterskip=3pt},
}

%% ---- 字体 ------------------------------------------------------------------
%% 西文优先用 TeX Gyre Termes 与 Heros。中文字体依次回退到 Fandol、Songti SC、SimSun、Noto。
\IfFontExistsTF{texgyretermes-regular.otf}{%
  \setmainfont{texgyretermes}[Extension=.otf,UprightFont=*-regular,BoldFont=*-bold,
    ItalicFont=*-italic,BoldItalicFont=*-bolditalic]%
  \setsansfont{texgyreheros}[Extension=.otf,UprightFont=*-regular,BoldFont=*-bold,
    ItalicFont=*-italic,BoldItalicFont=*-bolditalic]%
}{}
\IfFontExistsTF{FandolSong-Regular.otf}{%
  \setCJKmainfont{FandolSong-Regular.otf}[BoldFont=FandolSong-Bold.otf,ItalicFont=FandolKai-Regular.otf]%
  \setCJKsansfont{FandolHei-Regular.otf}[BoldFont=FandolHei-Bold.otf]%
  \setCJKmonofont{FandolFang-Regular.otf}%
}{%
  \IfFontExistsTF{Songti SC}{%
    \setCJKmainfont{Songti SC}[BoldFont=Heiti SC,ItalicFont=Kaiti SC]%
    \setCJKsansfont{Heiti SC}%
  }{%
    \IfFontExistsTF{SimSun}{%
      \setCJKmainfont{SimSun}[BoldFont=SimHei,ItalicFont=KaiTi]\setCJKsansfont{SimHei}%
    }{%
      \setCJKmainfont{Noto Serif CJK SC}\setCJKsansfont{Noto Sans CJK SC}%
    }%
  }%
}
%% 审稿原文里的希腊字母用数学字体显示，源文本保持原字符，缺字时在此增补。
\newunicodechar{α}{\ensuremath{\alpha}}
\newunicodechar{β}{\ensuremath{\beta}}
\newunicodechar{γ}{\ensuremath{\gamma}}
\newunicodechar{δ}{\ensuremath{\delta}}
\newunicodechar{ε}{\ensuremath{\varepsilon}}
\newunicodechar{λ}{\ensuremath{\lambda}}
\newunicodechar{μ}{\ensuremath{\mu}}
\newunicodechar{ρ}{\ensuremath{\rho}}
\newunicodechar{σ}{\ensuremath{\sigma}}
\newunicodechar{τ}{\ensuremath{\tau}}

%% ---- 颜色 ------------------------------------------------------------------
\definecolor{rbgray}{HTML}{4D4D4D}
\colorlet{rbcur}{rbgray}

%% ---- 审稿人信息与字符统计 -------------------------------------------------
\ExplSyntaxOn
\prop_new:N \g_rb_meta_prop
\prop_new:N \g_rb_en_prop
\prop_new:N \g_rb_zh_prop
\prop_new:N \g_rb_label_prop
\prop_new:N \g_rb_count_prop
\seq_new:N \g_rb_reviewers_seq
\seq_new:N \g_rb_fields_seq
\tl_new:N \l_rb_rows_tl
\tl_new:N \l_rb_head_tl
\tl_new:N \g_rb_cursec_tl   % 当前节（general 或审稿人编号）
\int_new:N \g_rb_part_int   % 当前节内的评论序号，由 \CommentBreak 递增
\int_new:N \l_rb_idx_int

%% \RevReviewer{编号}{颜色HTML} 写在 reviewers.tex 中，回复工程与修订稿共用这一份。
\NewDocumentCommand \RevReviewer { m m }
  {
    \seq_gput_right:Nn \g_rb_reviewers_seq {#1}
    \definecolor{rb@#1}{HTML}{#2}
  }
%% \SetScoreFields{Rating, Confidence, ...} 声明会议评分表的栏目，按审稿表填写。
\NewDocumentCommand \SetScoreFields { m }
  { \seq_gset_from_clist:Nn \g_rb_fields_seq {#1} }
%% \DefScores{编号}{值1, 值2, ...} 的顺序与 \SetScoreFields 相同。值里有逗号时用花括号括起来。
\NewDocumentCommand \DefScores { m m }
  {
    \int_zero:N \l_rb_idx_int
    \clist_map_inline:nn {#2}
      {
        \int_incr:N \l_rb_idx_int
        \prop_gput:Nxn \g_rb_meta_prop { #1/\int_use:N \l_rb_idx_int } {##1}
      }
  }

%% \DefConcern{意见编号}{显示标签}{英文原文}{中文译文}，见 reviews_bilingual.tex
\NewDocumentCommand \DefConcern { m m +m +m }
  {
    \prop_gput:Nnn \g_rb_label_prop {#1} {#2}
    \prop_gput:Nnn \g_rb_en_prop {#1} {#3}
    \prop_gput:Nnn \g_rb_zh_prop {#1} {#4}
  }

%% counts.tex 中的 \RBsetcount{文件名}{字符数}{上限}{未回填项}{空回复块}。
%% 文件名是节编号（如 R2），节内用了 \CommentBreak 时是 <节>_part<k>（如 R1_part2）。
%% \RBsetparts{节}{评论条数} 只对用了 \CommentBreak 的节写出。
\NewDocumentCommand \RBsetparts { m m }
  { \prop_gput:Nnn \g_rb_count_prop {#1/parts} {#2} }
\NewDocumentCommand \RBsetcount { m m m m m }
  {
    \prop_gput:Nnn \g_rb_count_prop {#1/n} {#2}
    \prop_gput:Nnn \g_rb_count_prop {#1/limit} {#3}
    \prop_gput:Nnn \g_rb_count_prop {#1/pending} {#4}
    \prop_gput:Nnn \g_rb_count_prop {#1/tbd} {#5}
  }
%% 只显示“字符数 / 上限”，超限标红
\NewDocumentCommand \RBcountshort { m }
  {
    \prop_if_in:NnTF \g_rb_count_prop {#1/n}
      {
        \int_compare:nNnTF { \prop_item:Nn \g_rb_count_prop {#1/n} }
            > { \prop_item:Nn \g_rb_count_prop {#1/limit} }
          { \textcolor{red}{\textbf{\prop_item:Nn \g_rb_count_prop {#1/n}}} }
          { \textbf{\prop_item:Nn \g_rb_count_prop {#1/n}} }
        \,/\,\prop_item:Nn \g_rb_count_prop {#1/limit}
      }
      { \textcolor{red}{未统计} }
  }
%% 完整统计行，#1 是行首文字，#2 是 counts.tex 中的文件名
\cs_new_protected:Npn \rb_count_line:nn #1#2
  {
    \prop_if_in:NnTF \g_rb_count_prop {#2/n}
      {
        #1~\RBcountshort{#2}
        \int_compare:nNnT { \prop_item:Nn \g_rb_count_prop {#2/n} }
            > { \prop_item:Nn \g_rb_count_prop {#2/limit} }
          { \textcolor{red}{（超出上限）} }
        \quad 未回填数字~
        \int_compare:nNnTF { \prop_item:Nn \g_rb_count_prop {#2/pending} } > { 0 }
          { \textcolor{red}{\prop_item:Nn \g_rb_count_prop {#2/pending}} }
          { 0 }
        ~处 \quad 待填写回复块~
        \int_compare:nNnTF { \prop_item:Nn \g_rb_count_prop {#2/tbd} } > { 0 }
          { \textcolor{red}{\prop_item:Nn \g_rb_count_prop {#2/tbd}} }
          { 0 }
        ~个
      }
      { \textcolor{red}{字符数未统计，先运行~make~check} }
  }
\NewDocumentCommand \RBcount { m } { \rb_count_line:nn {英文回复字符数} {#1} }
%% 用了 \CommentBreak 的节，第 k 条评论的统计行，#1 节，#2 序号
\cs_new_protected:Npn \rb_part_line:nn #1#2
  {
    \prop_if_in:NnTF \g_rb_count_prop {#1/parts}
      {
        第~#2~条评论（共~\prop_item:Nn \g_rb_count_prop {#1/parts}~条，导出为~#1\textunderscore{}part#2.md）\quad
        \rb_count_line:nn {本条} {#1_part#2}
      }
      { \textcolor{red}{counts.tex~中没有本节的分条统计，先运行~make~check} }
  }
%% 节标题下的统计行。分条的节显示第 1 条，不分条的节显示全节。
\NewDocumentCommand \RBsectioncount { m }
  {
    \prop_if_in:NnTF \g_rb_count_prop {#1/parts}
      { \rb_part_line:nn {#1} {1} }
      { \RBcount{#1} }
  }
%% 评分一览表的字符数格。分条的节逐条列出。
\NewDocumentCommand \RBcountcell { m }
  {
    \prop_if_in:NnTF \g_rb_count_prop {#1/parts}
      {
        \int_step_inline:nn { \prop_item:Nn \g_rb_count_prop {#1/parts} }
          {
            \int_compare:nNnT {##1} > {1} { ，}
            第~##1~条~\RBcountshort{#1_part##1}
          }
      }
      { \RBcountshort{#1} }
  }

%% ---- 意见显示 ----------------------------------------------------------------
\newtcolorbox{rbconcernbox}{enhanced,sidebyside,sidebyside~align=top~seam,
  sidebyside~gap=7pt,lefthand~ratio=0.5,colback=rbcur!4,colframe=rbcur!40,
  boxrule=0.4pt,arc=1pt,left=4pt,right=4pt,top=2pt,bottom=2pt,
  before~skip=3pt,after~skip=3pt,fontupper=\small,fontlower=\small,
  segmentation~style={solid,rbcur!30},
  before~upper={\setlength{\parskip}{3pt}},before~lower={\setlength{\parskip}{3pt}}}
\NewDocumentCommand \rbtag { m } { {\sffamily\bfseries\footnotesize\color{rbcur}#1}\enspace }

\NewDocumentCommand \ShowConcern { m }
  {
    \prop_if_in:NnTF \g_rb_en_prop {#1}
      {
        \begin{rbconcernbox}
          \rbtag{\prop_item:Nn \g_rb_label_prop {#1}~原文}
          \prop_item:Nn \g_rb_en_prop {#1}
          \tcblower
          \rbtag{\prop_item:Nn \g_rb_label_prop {#1}~译文}
          \prop_item:Nn \g_rb_zh_prop {#1}
        \end{rbconcernbox}
      }
      { \par\textcolor{red}{\textbf{未定义的意见编号~#1}}\par }
  }

%% rbgroup  一组意见加一个合并回复。
%% \begin{rbgroup}[可选标题]{编号1, 编号2, ...} 先排出所列意见的原文与译文，
%% 环境内再写 responseen 与 responsezh 两个回复环境。编号列表可以为空。
%% 带星号的 rbgroup* 只在组标题中列出编号，不重复排出原文与译文（用于总回复）。
\NewDocumentEnvironment { rbgroup } { s o m }
  {
    \par\addvspace{5pt}
    \noindent{\color{rbcur}\rule[0.5ex]{1.2em}{0.8pt}}\enspace
    {\sffamily\bfseries\small\color{rbcur}
      \IfValueT{#2}{#2\quad}
      \tl_if_blank:nF {#3} { 意见~\clist_use:nn {#3} {,~} } }
    \par\nobreak
    \IfBooleanF {#1} { \clist_map_inline:nn {#3} { \ShowConcern{##1} } }
  }
  { \par\addvspace{4pt} }

%% ---- 回复环境 ----------------------------------------------------------------
%% responseen 的内容由 tools/export_openreview.py 导出到 OpenReview，
%% 只能使用导出脚本支持的命令（见 README.md）。responsezh 不导出。
\tcbset{rbresp/.style={enhanced,breakable,boxrule=0.5pt,arc=1pt,
  left=5pt,right=5pt,top=3pt,bottom=3pt,before~skip=3pt,after~skip=3pt,
  fonttitle=\sffamily\bfseries\footnotesize,toptitle=0.6pt,bottomtitle=0.6pt,
  fontupper=\small,before~upper={\setlength{\parskip}{3pt}}}}
\newtcolorbox{responseen}{rbresp,colframe=rbcur,colbacktitle=rbcur,coltitle=white,
  colback=white,title={我方回复~英文~（导出到~OpenReview）}}
\newtcolorbox{responsezh}{rbresp,colframe=rbcur!45,colbacktitle=rbcur!12,coltitle=rbcur!80!black,
  colback=rbcur!2,title={我方回复~中文对照~（不导出）}}

%% 回复中可用的命令
\NewDocumentCommand \pending { m } { \textcolor{red}{\textbf{[待回填~#1]}} }
\NewDocumentCommand \respTBD { } { {\color{gray}\itshape（待填写）} }
\NewDocumentCommand \rhead { m } { \par\addvspace{2pt}\noindent\textbf{#1}\par\nobreak }
%% rtable  简单表格，导出为 Markdown 管道表。列格式只用 l c r，第一行是表头。
%% 列格式为 lll 的表（每位审稿人第 1 条的一览表）在 PDF 中按行宽折行，三列宽度比为 0.62 : 1.38 : 1，
%% 并在 \par 之后加 \nopagebreak[4]，使表前一段（表题）的末行与表格排在同一页。
%% 其余列格式整表缩放到行宽。导出的 Markdown 不受影响。
\NewDocumentEnvironment { rtable } { m +b }
  {
    \par
    \str_if_eq:nnT {#1} { lll } { \nopagebreak[4] }
    \addvspace{3pt}\begingroup\centering\small
    \str_if_eq:nnTF {#1} { lll }
      { \rb_rtable_wrap:n {#2} }
      { \begin{adjustbox}{max~width=\linewidth}\begin{tabular}{#1}\toprule #2 \crcr\bottomrule\end{tabular}\end{adjustbox} }
    \par\endgroup\addvspace{3pt}
  }
  { }
\cs_new_protected:Npn \rb_rtable_wrap:n #1
  {
    \begin{tabularx} { \linewidth }
      { @{} >{\raggedright\arraybackslash\hsize=0.62\hsize}X
            >{\raggedright\arraybackslash\hsize=1.38\hsize}X
            >{\raggedright\arraybackslash}X @{} }
    \toprule #1 \crcr\bottomrule
    \end{tabularx}
  }

%% ---- 节 ------------------------------------------------------------------------
\NewDocumentCommand \GeneralSection { }
  {
    \colorlet{rbcur}{rbgray}
    \tl_gset:Nn \g_rb_cursec_tl {general}
    \int_gset:Nn \g_rb_part_int {1}
    \section*{\color{rbgray}General~Response~（全体审稿人可见的总回复）}
    \pdfbookmark[1]{General~Response}{rb-general}
    {\small \RBsectioncount{general}}\par
  }
\NewDocumentCommand \ReviewerSection { m }
  {
    \clearpage
    \colorlet{rbcur}{rb@#1}
    \tl_gset:Nn \g_rb_cursec_tl {#1}
    \int_gset:Nn \g_rb_part_int {1}
    \section*{\color{rb@#1}审稿人~#1}
    \pdfbookmark[1]{Reviewer~#1}{rb-#1}
    {\small
      \seq_map_indexed_inline:Nn \g_rb_fields_seq
        { ##2~\prop_item:Nn \g_rb_meta_prop { #1/##1 } \quad }
      \par
      \RBsectioncount{#1}\par}
  }

%% \CommentBreak{标签}  在两个 rbgroup 之间开始同一节的下一条 OpenReview 评论。
%% 导出脚本把其后的回复块写入下一份文件（<节>_part2.md 等），每条各自统计字符数。
%% 标签写成 (2 of 3) 时，导出脚本核对它与实际条数一致。
\newtcolorbox{rbbreakbox}{enhanced,colback=rbcur!8,colframe=rbcur,boxrule=0pt,
  leftrule=3pt,arc=0pt,left=6pt,right=4pt,top=3pt,bottom=3pt,
  before~skip=10pt,after~skip=6pt,fontupper=\small}
\NewDocumentCommand \CommentBreak { m }
  {
    \int_gincr:N \g_rb_part_int
    \par
    \begin{rbbreakbox}
      {\sffamily\bfseries\color{rbcur}
        OpenReview~评论分界\quad #1}\par
      \use:x { \exp_not:N \rb_part_line:nn { \g_rb_cursec_tl } { \int_use:N \g_rb_part_int } }
    \end{rbbreakbox}
  }

%% 评分一览表，栏目由 \SetScoreFields 决定
\cs_new_protected:Npn \rb_add_row:n #1
  {
    \tl_put_right:Nn \l_rb_rows_tl { {\color{rb@#1}\bfseries #1} }
    \seq_map_indexed_inline:Nn \g_rb_fields_seq
      { \tl_put_right:Nn \l_rb_rows_tl { & \prop_item:Nn \g_rb_meta_prop { #1/##1 } } }
    \tl_put_right:Nn \l_rb_rows_tl { & \RBcountcell{#1} \\ }
  }
\NewDocumentCommand \RBScoreTable { }
  {
    \tl_clear:N \l_rb_rows_tl
    \tl_set:Nn \l_rb_head_tl { 审稿人 }
    \seq_map_inline:Nn \g_rb_fields_seq
      { \tl_put_right:Nn \l_rb_head_tl { & ##1 } }
    \seq_map_function:NN \g_rb_reviewers_seq \rb_add_row:n
    \tl_put_right:Nx \l_rb_rows_tl
      {
        \exp_not:n { \midrule 总回复 }
        \prg_replicate:nn { \seq_count:N \g_rb_fields_seq } { \exp_not:n { & } }
        \exp_not:n { & \RBcountcell{general} \\ }
      }
    \begin{center}\small
    \use:x
      {
        \exp_not:N \begin{tabular}
          { l \prg_replicate:nn { \seq_count:N \g_rb_fields_seq } { c }
            > { \exp_not:N \raggedright \exp_not:N \arraybackslash } p { 5.3cm } }
      }
      \toprule
      \l_rb_head_tl & 英文回复字符数 \\
      \midrule
      \l_rb_rows_tl
      \bottomrule
    \end{tabular}
    \end{center}
  }

%% 附录中的参考意见（灰色，不回复）
\NewDocumentCommand \ReferenceSection { m }
  {
    \clearpage
    \colorlet{rbcur}{rbgray}
    \section*{\color{rbgray}#1}
    \pdfbookmark[1]{#1}{rb-ref}
  }
\ExplSyntaxOff
